% Paper 2: preference-conditioned policy learning for the four-objective lot-sizing model,
% trained by reinforcement learning and evaluated against exact Pareto fronts.
% Target: Computers & Operations Research (subscription route, no publication fee);
% alternative: International Journal of Production Economics.
% Build: python3 make_tables.py --run ../ --out .   then   pdflatex main && bibtex main && pdflatex main && pdflatex main
\documentclass[preprint,12pt,authoryear]{elsarticle}
\usepackage[T1]{fontenc}
\usepackage{lmodern}
\usepackage{amsmath,amssymb,amsthm,mathtools}
\usepackage{booktabs}
\usepackage{adjustbox}
\usepackage{graphicx}
\usepackage[hidelinks]{hyperref}
\usepackage[capitalise,noabbrev]{cleveref}
\newtheorem{proposition}{Proposition}
\newtheorem{definition}{Definition}
\newcommand{\tabin}[1]{\IfFileExists{tables/#1}{\begin{adjustbox}{max width=\linewidth}\input{tables/#1}\end{adjustbox}}{\fbox{run make\_tables.py: \detokenize{#1}}}}
\newcommand{\figin}[2]{\IfFileExists{figures/#2}{\includegraphics[width=#1]{figures/#2}}{\fbox{run make\_tables.py: \detokenize{#2}}}}
\IfFileExists{tables/results_macros.tex}{% generated by make_tables.py
\makeatletter
\@namedef{rv@SynN}{300}
\@namedef{rv@SynPairs}{25200}
\@namedef{rv@RealN}{700}
\@namedef{rv@SynMedDec}{0.0215}
\@namedef{rv@SynMedBest}{0.0107}
\@namedef{rv@SynMedGreedy}{0.0375}
\@namedef{rv@RealMedDec}{0.0318}
\@namedef{rv@RealMedBest}{0.0153}
\@namedef{rv@RealMedGreedy}{0.0438}
\@namedef{rv@RealExactMed}{0.20}
\@namedef{rv@RealDecMs}{13.0}
\@namedef{rv@RealPolMs}{214}
\@namedef{rv@SynExactMax}{259}
\@namedef{rv@SynExactP}{42}
\@namedef{rv@AblSpeed}{1.44}
\@namedef{rv@TrainHours}{2.3}
\@namedef{rv@MPL}{3.11.2}
\@namedef{rv@CPU}{12th Gen Intel Core i5-1235U}
\@namedef{rv@GCC}{15.2.0}
\@namedef{rv@Python}{3.14.4}
\@namedef{rv@AblSentence}{see the ablation section.}
\newcommand{\Rv}[1]{\@ifundefined{rv@#1}{??}{\@nameuse{rv@#1}}}
\makeatother
}{\newcommand{\Rv}[1]{??}}
\newcommand{\POC}{\mathrm{POC}}\newcommand{\HC}{\mathrm{HC}}\newcommand{\CSL}{\mathrm{CSL}}\newcommand{\ITO}{\mathrm{ITO}}

\journal{Computers \& Operations Research}

\begin{document}
\begin{frontmatter}
\title{Learning to plan lot sizes for any trade-off: a preference-conditioned policy on a provably
sufficient action space, measured against exact Pareto fronts}

\author[unimed]{Michael Sipayung}
\address[unimed]{Department of Physics, Universitas Negeri Medan, Medan, Indonesia}

\begin{abstract}
Decision makers in lot sizing weigh ordering cost, holding cost, service level and inventory
turnover against one another, and their preferred trade-off changes from item to item and from
week to week. We train a single neural policy that, given an instance and a preference vector,
constructs a lot-sizing plan for the four-objective model of Kania et al.\ in milliseconds. Two
features distinguish the approach from existing learning-based methods for inventory and for
multiobjective combinatorial optimisation. First, the policy acts on order events rather than on
per-period quantities: by an order-event theorem for this model, the event action space contains
every Pareto-optimal plan while being much smaller than the space of quantity sequences. Second,
every result is measured against the exact Pareto front, computed by the label-setting dynamic
programme of our companion paper, so that the reported gaps are true optimality gaps rather than
gaps to a heuristic reference. On \Rv{SynN} held-out synthetic instances the policy attains a
median Tchebycheff regret of \Rv{SynMedDec} (\Rv{SynMedBest} with best-of-17 sampling) against
\Rv{SynMedGreedy} for a greedy constructive heuristic, and on \Rv{RealN} instances built from the
M5 and Online Retail II data, which it never saw in training, \Rv{RealMedDec}
(\Rv{RealMedBest}) against \Rv{RealMedGreedy}. Learning on order events is \Rv{AblSpeed} times
faster to train than learning on per-period quantities and gives lower median and 90th-percentile
regret at equal iterations and at equal training time, while the quantity formulation lands
exactly on the front more often, mainly at extreme preferences.
We also report where the exact method remains preferable: for small real instances the full
front is computed faster than the policy answers a single preference with sampling.
\end{abstract}

\begin{keyword}
Inventory \sep lot sizing \sep multiobjective optimisation \sep Pareto set learning \sep
reinforcement learning \sep exact benchmark
\end{keyword}
\end{frontmatter}

% =====================================================================================
\section{Introduction}\label{sec:intro}

Lot-sizing decisions in practice are rarely made against a single cost. Kania and co-authors
formulate the single-item problem with safety stock and safety order time as a four-objective
model, in which purchasing and ordering cost, holding cost, cycle service level and inventory
turnover are traded off by a decision maker who explores the Pareto front interactively
\citep{Kania2021ICCL,Kania2022CIE,Kania2024DESMILS}. Interactive exploration, however, requires
that a plan matching the current preference can be produced quickly, again and again, for every
item that the decision maker is responsible for. Our companion paper makes the exact Pareto front
of this model computable by a label-setting dynamic programme; the front is exact, but its size
grows to tens or hundreds of thousands of points and its computation, while fast for most real
series, takes minutes on the hardest instances.

This paper asks whether a single learned policy can answer any preference directly, and how far
its answers are from optimal. The question belongs to Pareto set learning (PSL), in which one
preference-conditioned model approximates the whole Pareto set \citep{Lin2022PSL}. PSL has been
developed mainly for routing and knapsack problems, where exact Pareto fronts are out of reach and
quality is therefore reported relative to heuristic references such as weighted-sum LKH or
hypervolume of an approximate front \citep{Chen2023NHDE,Lu2024GAPL}. In inventory control, deep
reinforcement learning has been shown to match or beat strong heuristics in single-objective
settings \citep{Gijsbrechts2022MSOM,Dehaybe2024,vanHezewijk2023,vanHezewijk2025}, and the few
multiobjective studies approximate fronts with evolutionary search over policy parameters, again
without an exact reference \citep{Kotecha2025MORSE,Rachman2025}. \citet{Boute2022Roadmap} identify
the use of structural policy insights from inventory theory as one of the main open avenues for
deep reinforcement learning in this field.

Our setting allows both gaps to be closed at once. The exact front gives every learned plan a true
regret, and a structural result for the model, the order-event theorem of the companion paper,
gives the policy an action space that is provably sufficient. Our contributions are:
\begin{enumerate}
\item An order-event Markov decision process for the four-objective model in which every
episode yields a feasible plan, every Pareto-optimal plan is reachable, and an episode has as many
decisions as the plan has orders (\cref{sec:mdp}). We verify the second property exhaustively on
small instances.
\item A preference-conditioned policy with a class head and a pointer-style event head, trained by
REINFORCE with a shared per-instance baseline on weighted Tchebycheff rewards normalised by the
exact ideal and nadir points (\cref{sec:policy}).
\item An evaluation protocol in which every (instance, preference) pair is scored by its regret
to the exact optimum of the same scalarisation, on held-out synthetic instances and on 700
real-demand instances (\cref{sec:design,sec:results}).
\item An ablation against the standard per-period quantity formulation at equal iterations and at
equal training time, isolating the value of the structural action space (\cref{sec:ablation}).
\item An honest account of when learning pays off and when the exact method should be used
instead (\cref{sec:discussion}).
\end{enumerate}

% =====================================================================================
\section{Related work}\label{sec:related}

\paragraph{Pareto set learning.}
\citet{Lin2022PSL} train a single attention model conditioned on a preference vector with a
multiobjective variant of POMO \citep{Kwon2020POMO}; the approach is a learning counterpart of
decomposition-based evolutionary search \citep{ZhangLi2007MOEAD}. Later work improves diversity
and geometry of the learned front \citep{Chen2023NHDE,Lu2024GAPL}. Because the benchmark problems
are NP-hard and large, these studies measure hypervolume relative to the best known approximate
front. Our problem is different in two respects: its exact front is computable, and its
constructive decisions can be restricted by theory without loss of optimality.

\paragraph{Reinforcement learning for lot sizing and inventory control.}
\citet{Gijsbrechts2022MSOM} and the roadmap of \citet{Boute2022Roadmap} establish deep
reinforcement learning as a credible tool for inventory problems. For lot sizing,
\citet{Dehaybe2024} train agents on randomly generated environments that generalise to
non-stationary demand, \citet{vanHezewijk2023,vanHezewijk2025} address the capacitated multi-item
case, and \citet{Felizardo2024} the discrete lot-sizing problem on parallel machines. All of these
optimise a single cost, and all act on per-period quantities. Multiobjective supply chain studies
\citep{Kotecha2025MORSE,Rachman2025} combine reinforcement learning with evolutionary search and
report hypervolume without an exact reference.

\paragraph{Exact multiobjective methods.}
The companion paper computes the exact four-objective front by a label-setting dynamic programme
in the tradition of multicriteria shortest paths \citep{Martins1984,SedenoNoda2019,Maristany2021};
we use it here as an oracle, not as a competitor, except in the timing comparison of
\cref{sec:discussion}.

% =====================================================================================
\section{Model and exact reference}\label{sec:model}

We use the model and notation of the companion paper. An instance is a demand forecast
$D_1,\dots,D_T$ with standard deviation $\sigma$, lead time $L$, order lattice step $r$ and
minimum order quantity, purchase price $p$, ordering cost $c$, holding rate $h$ and service
target. A plan fixes a safety strategy, summarised by its requirement class (the service level
$\CSL=\Phi(B/\sigma)$ is determined by the class), and an order schedule; the four objectives are
$\POC$ and $\HC$ (minimised) and $\CSL$ and $\ITO$ (maximised). The exact front $\mathcal{F}$ of
an instance, its ideal point $z^{\mathrm{I}}$ and nadir point $z^{\mathrm{N}}$ are computed once
per instance by the companion algorithm.

\paragraph{Scalarisation and regret.}
For a preference $\lambda$ in the unit simplex and the minimisation vector
$f=(\POC,\HC,-\CSL,-\ITO)$ we use the normalised weighted Tchebycheff function
\begin{equation}
g(f\mid\lambda)=\max_{i}\ \lambda_i\,\frac{f_i-z^{\mathrm{I}}_i}{\rho_i},\qquad
\rho_i=\max\bigl(z^{\mathrm{N}}_i-z^{\mathrm{I}}_i,\ 10^{-3}s_i\bigr),
\end{equation}
where $s_i$ is a natural scale of objective $i$ for the instance (purchase value of total demand
plus ordering cost for $\POC$, holding cost of total demand over the horizon for $\HC$, one for
$\CSL$, and a turnover scale for $\ITO$). The floor matters only when an objective is constant on
the front, e.g.\ when every Pareto plan has $\CSL=1$ or when initial stock covers the whole
horizon so that $\POC=0$ on the front. The regret of a plan with objective vector $f$ is
\begin{equation}
\mathrm{reg}(f\mid\lambda)=g(f\mid\lambda)-\min_{q\in\mathcal{F}}g(q\mid\lambda)\ \ge 0,
\end{equation}
and is zero exactly when the plan is optimal for the scalarisation. Since every
Pareto-optimal point minimises $g$ for some $\lambda$, regret over a grid of preferences measures
both convergence and coverage.

% =====================================================================================
\section{An order-event Markov decision process}\label{sec:mdp}

The order-event theorem of the companion paper states that, within a requirement class, an order
placed in period $\tau$ before the next order in period $\tau'$ can be restricted to the least
feasible lattice level $X_0=\max(\lceil N\rceil_r, x'+\mathrm{moq})$ together with a short list of
exceptional levels that are needed only while the projected stock at $\tau'$ would fall below the
positivity threshold. We turn this into a Markov decision process.

\paragraph{States and actions.}
The first decision selects a requirement class. Thereafter a state is
$(\tau,x,k,W,R)$: the current order period, the cumulative quantity before the order, the number
of orders so far, and the accumulated holding and turnover terms. An action is a pair
$(\tau',X)$ consisting of the next order period and the level reached by the current order, drawn
from the levels permitted by the theorem; $\tau'=T'+1$ ends the episode. Before the first order,
actions choose the length of the order-free prefix. Objective increments are additive, so the
final objective vector is available when the episode ends.

\begin{proposition}\label{prop:mdp}
(i) Every episode of the order-event process ends in a feasible plan. (ii) Every Pareto-optimal
plan of the instance is the outcome of some episode. (iii) An episode has one decision per order
plus two.
\end{proposition}

Part (i) holds because every action satisfies the requirement, positivity and lattice constraints
of the periods it covers; (iii) is immediate. Part (ii) follows from the order-event theorem and
the class-dominance corollary of the companion paper. We also verified (ii) computationally: on
small instances, exhaustive enumeration of all episodes reproduces the exact front point for point
(\cref{tab:check}).

\begin{table}[t]\centering
\caption{Exhaustive check of \cref{prop:mdp}(ii): all episodes of the order-event process were
enumerated and filtered for non-dominance.}\label{tab:check}
\tabin{check.tex}
\end{table}

\paragraph{The per-period alternative.}
The usual formulation in reinforcement learning for lot sizing lets the agent choose, in every
period, either no order or an order quantity from a discretised set
\citep{Gijsbrechts2022MSOM,Dehaybe2024,vanHezewijk2023}. For the ablation we implement this
formulation on the same model: in period $u$ the agent may skip (if the current level remains
feasible) or raise the cumulative quantity to one of 24 lattice levels starting from the least
feasible one; infeasible actions are masked, so every episode is again a feasible plan, mapped
onto the same objectives. Episodes then have $T'$ decisions regardless of the number of orders.

% =====================================================================================
\section{Policy and training}\label{sec:policy}

\paragraph{Architecture.}
The policy has two heads, each a multilayer perceptron applied to every candidate and followed by
a softmax over candidates, in the manner of pointer networks. The class head scores requirement
classes from nine features (normalised service buffer, service level, mean and maximum
requirement, share of periods with positive requirement, and the preference). The event head
scores each candidate $(\tau',X)$ from 23 features describing the state (time, cumulative
quantity, orders, accumulated terms), the candidate (horizon covered, level, order size, cost
increments, slack to the requirement, terminal flag) and the preference. All quantities are
normalised by instance scales, so that one network serves instances of very different magnitude.
The heads have two hidden layers of 64 and 96 units with $\tanh$ activations.

\paragraph{Training.}
For a batch of (instance, preference) pairs, with preferences drawn uniformly from the simplex,
we sample $K=8$ episodes per pair and use the reward $-g(f\mid\lambda)$. Following POMO
\citep{Kwon2020POMO,Lin2022PSL}, the baseline is the mean reward of the $K$ episodes of the same
pair, and advantages are standardised within the group; the gradient is the REINFORCE estimator
\citep{Williams1992}, applied with Adam \citep{KingmaBa2015} at learning rate $10^{-3}$ and
gradient-norm clipping. The whole implementation, including backpropagation, is written in C++
with OpenMP and verified against finite differences (maximum deviation $\approx10^{-10}$).

\paragraph{Training distribution.}
Training instances are generated with $T=48$, mean demand log-uniform on $[1,200]$, coefficient
of variation in $[0.1,1]$, lead time 1--5, lattice step and minimum order tied to mean demand,
price log-uniform, holding rate proportional to price, ordering cost scattered around the EOQ
balance, and three demand regimes: stationary (60\%), intermittent (20\%) and seasonal (20\%).
The pool contains 2000 training and 300 held-out instances, each with its exact front.

% =====================================================================================
\section{Experimental design}\label{sec:design}

\paragraph{Instances.}
Held-out synthetic instances come from the same generator with disjoint seeds. Real instances are
the 700 benchmark instances of the companion paper, built from the M5 \citep{Makridakis2022M5}
and Online Retail II \citep{Chen2012Retail} data across the Syntetos--Boylan demand classes
\citep{SyntetosBoylan2005}; their exact fronts are taken from the companion runs. The policy never
sees real data during training.

\paragraph{Preferences and metrics.}
Each instance is evaluated on the 84 preferences of the simplex lattice with step $1/6$. For every
pair we record the regret of the greedy decoding, of the best of greedy decoding and 16 sampled
episodes, and of the baseline below; we report median, mean, 90th percentile and maximum, and the
share of pairs whose plan lies on the exact front within the tolerances of the companion paper.
Inference times are measured on one thread.

\paragraph{Baseline.}
A greedy constructive heuristic enumerates all requirement classes and, within each, repeatedly
chooses the candidate whose cost increment per covered period, weighted by the preference and
normalised by the front ranges, is smallest; the best class under $g$ is returned. It uses the exact ideal and nadir points,
which favours it.

\paragraph{Hardware and software.}
All runs used one laptop (\Rv{CPU}, 9 threads) and GCC \Rv{GCC} with options \texttt{-O3} and \texttt{-march=native} and
OpenMP. Training of the main policy took \Rv{TrainHours} hours.

% =====================================================================================
\section{Results}\label{sec:results}

\subsection{Held-out synthetic instances}
\begin{table}[t]\centering
\caption{Regret to the exact optimum on \Rv{SynN} held-out synthetic instances
($\Rv{SynPairs}$ instance--preference pairs).}\label{tab:syn}
\tabin{syn.tex}
\end{table}

\Cref{tab:syn} reports the regret on the held-out synthetic instances. With greedy decoding the
policy attains a median regret of 0.0215 and a 90th percentile of 0.096, against 0.0375 and 0.282
for the greedy heuristic; best-of-17 sampling halves the median to 0.0107. The policy is within
0.05 of the optimum on 75.6\% of the pairs (85.1\% with sampling) against 55.6\% for the
heuristic, and it fails badly (regret above 0.5) on 0.39\% of the pairs against 2.28\%
(\cref{tab:dist}). The heuristic, however, is exactly optimal more often (11.6\% against 5.9\%):
enumerating every class and choosing the cheapest increments recovers the extreme plans, while the
policy is usually close to, but not exactly at, the optimum. Its maximum regret (1.38) also exceeds
that of the heuristic (0.78); the worst cases are concentrated on a few instances with a narrow
front in one objective, where the normalisation magnifies small deviations.

\begin{table}[t]\centering
\caption{Distribution of regret (shares of instance--preference pairs).}\label{tab:dist}
\tabin{dist.tex}
\end{table}



\subsection{Real demand: M5 and Online Retail II}
\begin{table}[t]\centering
\caption{Regret to the exact optimum on the \Rv{RealN} real-demand instances of the companion
paper; the policy was trained on synthetic instances only.}\label{tab:real}
\tabin{real.tex}
\end{table}

On the 700 real instances (\cref{tab:real}) the policy, trained on synthetic data only, keeps
its advantage: median regret 0.0318 against 0.0438 for the heuristic (0.0153 with sampling),
90th percentile 0.219 against 0.338, and 62.4\% of pairs within 0.05 of the optimum against
52.5\%. The share of severe failures is the same for both (4.75\% and 4.73\%) and higher than on
synthetic data; it comes from real series whose fronts are very narrow in service level, so that
choosing a lower service class costs hundreds of front ranges. For this reason we report medians,
percentiles and failure shares rather than means for the real data. The gain is larger on the M5
series (median 0.0302) than on the Online Retail series (0.0445), whose lumpy demand is less well
represented by the training generator.

\subsection{Action-space ablation}\label{sec:ablation}
\begin{table}[t]\centering
\caption{Order events versus per-period quantities: same pool, network, batch and optimiser;
compared at equal iterations (5000) and at equal training time.}\label{tab:abl}
\tabin{ablation.tex}
\end{table}

\Cref{tab:abl} compares the two action spaces. Training on per-period quantities takes 1.44 times
longer for the same 5000 iterations, and greedy decoding is 1.3--1.6 times slower at inference
(26.1 against 19.8\,ms per preference on synthetic instances, 21.4 against 13.0\,ms on real ones),
because an episode has one decision per period instead of one per order. At equal iterations the
order-event policy has the lower median (0.0215 against 0.0262 on synthetic, 0.0318 against
0.0353 on real data), the lower 90th percentile (0.096 against 0.127; 0.219 against 0.284), fewer
severe failures (0.39\% against 1.10\%; 4.75\% against 6.85\%) and a three times smaller worst case
on synthetic data. At equal training time the gap widens (median 0.0352 and 0.0457 for the
quantity policy).

The quantity policy nevertheless reaches the exact front more often (23--29\% of pairs against
10--12\%), and \cref{tab:pref} shows where: at the vertices of the preference simplex its median
regret is zero, i.e.\ it learns to produce the extreme plans exactly, whereas on edge and interior
preferences, which make up 95\% of the grid, the order-event policy is better. The two action
spaces therefore differ not only in how large their errors are but in where they make them. On the
100 Online Retail instances the quantity policy is better also in median (0.0292 against 0.0445),
and with best-of-17 sampling its share of severe failures on real data is slightly lower (2.29\%
against 2.87\%). We read these results as a consistent advantage of the order-event space in
typical and tail regret and in computational cost, not as dominance in every statistic.

\begin{table}[t]\centering
\caption{Median regret (greedy decoding) by preference type: vertex (one non-zero weight), edge
(two) and interior (three or four).}\label{tab:pref}
\tabin{preftype.tex}
\end{table}

\begin{figure}[t]\centering
\figin{0.8\textwidth}{fig_learning.pdf}
\caption{Training curves of the two action spaces against wall-clock time.}\label{fig:learning}
\end{figure}

% =====================================================================================
\section{When does learning pay off?}\label{sec:discussion}

The exact algorithm returns the whole front at once; the policy returns one plan per preference.
On the real instances the median exact computation takes \Rv{RealExactMed}\,s for the full front,
while the policy needs \Rv{RealDecMs}\,ms for one preference with greedy decoding and
\Rv{RealPolMs}\,ms with 16 samples. For a decision maker exploring a few dozen preferences on a
typical real series, the exact method is therefore both faster and optimal. The policy pays off
on hard instances, where exact computation takes up to \Rv{SynExactMax}\,s (90th percentile
\Rv{SynExactP}\,s on the synthetic pool), on large item portfolios that are re-planned frequently,
and as a component of models for which exact fronts are no longer computable, such as multi-item
problems with shared capacity \citep{Kania2024DESMILS}. In all cases the exact method remains the
reference that certifies the quality of the learned plans.

% =====================================================================================
\section{Conclusion}\label{sec:conclusion}
We trained a single preference-conditioned policy for the four-objective lot-sizing model and
measured it, pair by pair, against exact Pareto fronts. Trained only on synthetic instances, it
answers any preference in about 13--20\,ms with a median regret of 0.02--0.03 in normalised
Tchebycheff units on held-out synthetic and on real M5 and Online Retail instances, and reduces
the 90th percentile of the regret of a greedy constructive heuristic by 35--66\%. Building the action space on the order-event theorem of the
companion paper makes learning faster and the policy more accurate in typical and tail regret than
the per-period formulation that is standard in reinforcement learning for inventory, which in turn
is better at reproducing extreme plans exactly. The exact method remains indispensable: it
certifies the regret of every learned plan and, for small real instances, is itself the faster
way to the whole front.

Three directions follow. The policy is usually near-optimal but rarely exactly optimal; combining
it with the exact algorithm, for instance by letting the policy select a few requirement classes
and solving the inner problem exactly for those classes only, may close the remaining gap at a
fraction of the cost of the full front. A training distribution closer to lumpy retail demand
should reduce the gap observed on the Online Retail series. Finally, the multi-item setting of
\citet{Kania2024DESMILS}, where exact fronts are out of reach, is the natural next domain, with
the single-item exact method serving as a component and a benchmark.


\bibliographystyle{elsarticle-harv}
\bibliography{refs}
\end{document}
